\documentclass[11pt]{article}
\usepackage[margin=1in]{geometry}
\usepackage{amsmath,amssymb,amsthm,array,booktabs,expl3,hyperref}
% Theorem environments, notation and the \lean, \leanok and \uses
% tags are shared
% with the Lean blueprint (src/lean/blueprint), whose content is the appendix.
% `just blueprint-check` verifies every \lean{...} name here against
% the library.
% Macros shared by the web blueprint, the printed blueprint and the manuscript
% (src/latex/projects/free-foundations/src/sections/), so a statement reads the same wherever it appears.
% Theorem-like environments below appear in the dependency graph.
\newtheorem{theorem}{Theorem}
\newtheorem{proposition}[theorem]{Proposition}
\newtheorem{lemma}[theorem]{Lemma}
\newtheorem{corollary}[theorem]{Corollary}

\theoremstyle{definition}
\newtheorem{definition}[theorem]{Definition}
\theoremstyle{plain}

\newcommand{\R}{\mathbb{R}}
\newcommand{\E}{\mathbb{E}}
\newcommand{\norm}[1]{\left\lVert#1\right\rVert}
\newcommand{\abs}[1]{\left\lvert#1\right\rvert}
\newcommand{\ip}[2]{\left\langle#1,#2\right\rangle}
\newcommand{\tr}{\operatorname{tr}}
\newcommand{\KL}{\operatorname{KL}}
\newcommand{\diag}{\operatorname{diag}}

% Print versions of the blueprint macros, used by print.tex and by the manuscript.
% plasTeX defines the web versions. In print, a bound statement opens with a line
% naming the Lean declarations it is checked against, ticked when \leanok follows.
\ExplSyntaxOn
\seq_new:N \l__bp_names_seq
\NewDocumentCommand{\lean}{m}
 {
  \seq_clear:N \l__bp_names_seq
  \clist_map_inline:nn {#1} { \seq_put_right:Nn \l__bp_names_seq { \nolinkurl{##1} } }
  \par
  \group_begin:
  \upshape\footnotesize\raggedright Lean:~\seq_use:Nn \l__bp_names_seq {,~}
  % Peek last, so the next token is the document's.
  \peek_meaning_remove_ignore_spaces:NTF \leanok
   { ~$\checkmark$ \par \group_end: }
   { \par \group_end: }
 }
% Labels in \uses and \proves must exist: invisible references warn if they do not.
\NewDocumentCommand{\uses}{m}
 {\clist_map_inline:nn{#1}{\vphantom{\ref{##1}}}\ignorespaces}
\NewDocumentCommand{\proves}{m}
 {\clist_map_inline:nn{#1}{\vphantom{\ref{##1}}}\ignorespaces}
\ExplSyntaxOff
% A \leanok not consumed by \lean marks a finished proof.
\newcommand{\leanok}{{\footnotesize$\checkmark$}\ \ignorespaces}
\newcommand{\mathlibok}{{\footnotesize[Mathlib]}\ \ignorespaces}
\newcommand{\notready}{}
\newcommand{\discussion}[1]{}

\title{Measuring Financial Proximity with Pretrained Tabular Representations}
\author{Marcus Gawronsky\\\small\texttt{gwrmar002@myuct.ac.za}}
\date{Working research design --- October 2026}
\begin{document}
\maketitle
\begin{abstract}
  Spatial models of financial dependence require a measure of
  economic proximity.
  We propose a comparison between characteristic-based networks and networks
  constructed from pretrained representations of regulatory disclosures.
  The design separates the information supplied to an encoder from the geometry
  it induces, uses disclosures available before the evaluation period, and holds
  the spatial specification fixed across networks. Four separate questions
  organize the argument: which economic state matters, what a training
  objective preserves, whether the state can be recovered from a representation
  under held-out validation, and what the resulting network changes in
  covariance. A conditional perturbation result connects network measurement
  error to covariance error. The empirical
  hypothesis is that learned proximity organizes subsequent return dependence
  beyond conventional factors and characteristic distances. This working draft
  specifies the study; it does not report confirmatory empirical findings.
\end{abstract}

\section{Question and contribution}
Fern\'andez~\cite{fernandez} constructs financial proximity from histories of
individual firm characteristics and estimates spatial CAPM specifications.
Our question is whether a pretrained representation of several fundamentals
provides a more informative network under a fixed spatial-error model.
The contribution must be an economically interpretable measurement improvement,
conditional on the available information, rather than a claim that an encoder
creates an asset-pricing theory. We retain the original characteristic networks,
label combinations of their distances as extensions, and compare learned and
classical representations supplied with the same inputs.

Finance Research Letters is the provisional outlet. The intended main text is
short and centered on one empirical question and one analytical bridge;
additional model variants, data mappings, and proofs belong in supplementary
material. The current journal formatting requirements must be checked against
the publisher's author guide before submission.

Closely related work already studies learned asset representations.
Gabaix et al.~\cite{assetembeddings} construct representations from portfolio
holdings and examine return comovement, among other outcomes. Our proposed
comparison instead uses regulatory disclosures and externally pretrained
encoders to measure a network in a fixed spatial specification. Novelty must
be established against that literature as well as the literature on estimating
spatial weights. The perturbation proposition below is a standard
resolvent-based argument supporting the design, not itself a claim of a new
learning theorem.
F\"ulle and Otto~\cite{fulleotto} also construct spatial proximity in a
covariate space for financial volatility models. Our empirical distinction
must therefore isolate externally pretrained representation geometry and its
subsequent information content, rather than claim that financial proximity
from covariates is new.

\section{Economic mechanism and geometry}
Let $S_i$ describe a firm's operating and financing state and let
$X_i=g(S_i,V_i)$ denote its disclosures. After removing conventional factors,
suppose residual returns contain exposure to omitted economic shocks:
\begin{equation}
  u_i=b(S_i)'q+\eta_i,\qquad \operatorname{Var}(q)=\Omega.
\end{equation}
If $\norm{\Omega^{1/2}(b(S_i)-b(S_j))}_2\leq Ld_S(S_i,S_j)$,
then proximity bounds the difference in common-shock exposures. This alone
does not imply larger covariance: loading magnitudes and idiosyncratic risk
also matter. It motivates tests controlling for conventional exposures and
industry, rather than a claim of causal shock transmission.

The tabular foundation models evaluated here are in-context predictors, so each
row state depends on a declared support table and label:
$z_i=E(X_i\mid X_s,y_s)$. We distinguish \emph{state} readouts taken before
in-context prediction from \emph{task-conditioned} readouts that encode the
declared support target, and report both; neither is a target-neutral firm
embedding by construction. With normalized vectors $v_i=z_i/\norm{z_i}_2$,
define $d_{ij}=\norm{v_i-v_j}_2$ and
\begin{equation}
  W_{ij}=\frac{\exp(-d_{ij}/h)}{\sum_{k\ne i}\exp(-d_{ik}/h)},\quad i\ne j,
  \qquad W_{ii}=0.
\end{equation}
The source-faithful comparison uses $h=1$, as in Fern\'andez. Because distance
scales differ across representations, two predeclared alternatives apply the
same rule to classical and learned distances: $h$ equal to the median positive
distance, and $h$ chosen so that the mean effective neighbor count
$N^{-1}\sum_i{(\sum_j W_{ij}^2)}^{-1}$ equals that of the joint characteristic
network in the same year. The latter separates neighbor ranking from the amount
of smoothing. No bandwidth is chosen using returns.

Prediction alone does not identify this geometry. For any invertible $T$,
$\widetilde E=T E$ and $\widetilde h=h_{\rm pred}\circ T^{-1}$ preserve a
predictor while potentially changing distances. For example, transforming
$(1,0)$ and $(1,1)$ by $T=\operatorname{diag}(c,1)$ makes their normalized
chord distance approach zero as $c$ grows. Orthogonal invariance of chord
distance does not resolve general invertible transformations.

\begin{proposition}[Conditional measurement stability]\label{prop:stability}
  \lean{FreeFoundations.Stability.matrix_covariance_sub_le,
  FreeFoundations.Kernel.expKernel_linfty_le}\leanok%
  \uses{thm:stability-covariance, thm:kernel-sensitivity}
  Consider the same firms and diagonal innovation covariance $D$ for two spatial
  networks $W$ and $W^\star$. Put
  $A={(I-\lambda W)}^{-1}$ and $A^\star={(I-\lambda W^\star)}^{-1}$.
  If $\norm{A}_2,\norm{A^\star}_2\leq M$, then
  \begin{equation}
    \norm{ADA'-A^\star D(A^\star)'}_2
    \leq 2\norm{D}_2 |\lambda|M^3\norm{W-W^\star}_2.
  \end{equation}
  If distances in the exponential kernels satisfy
  $\max_{i\ne j}|d_{ij}-d^\star_{ij}|\leq\epsilon$, with common bandwidth $h$,
  then the maximum absolute row-sum norm obeys
  \begin{equation}
    \norm{W-W^\star}_\infty
    \leq \min\{2,\exp(2\epsilon/h)-1\}.
  \end{equation}
\end{proposition}

This is a conditional robustness result, not a theorem that pretraining
recovers $S_i$. The resolvent bound is an explicit assumption; row normalization
and $|\lambda|<1$ do not imply $\norm{A}_2\leq{(1-|\lambda|)}^{-1}$ for a
nonsymmetric network. Converting the row-norm bound to a spectral-norm bound
can introduce a factor $\sqrt{N}$, so no dimension-free conclusion is claimed.
The stronger learning-to-economic-distance link remains an empirical question
or requires additional identifiable economic supervision.

\section{What a learned measurement must establish}
A pretrained representation enters the spatial model only through distances.
Whether those distances measure economic proximity splits into four questions,
each with its own object and its own evidence; a positive answer to one does
not carry over to the next (Table~\ref{tab:framework}). The machine-checked
counterparts are in the accompanying Lean library; Appendix~\ref{sec:bp-reading}
restates each of its results, and every statement tagged ``Lean'' names the
declaration it is checked against.

\begin{table}[ht]
  \centering\small
  \begin{tabular}{*{4}{>{\raggedright\arraybackslash}p{0.22\linewidth}}}
    \toprule
    Question & Object & Formal link & Evidence required \\
    \midrule
    Which economic distinction matters? & $s_{\mathcal C}(H)$ & exposure
    geometry (\texttt{Exposure}) & a declared target and its reliability \\
    What does the objective preserve? & $F_m(H)$ & proper scores,
    sufficiency (\texttt{Score}, \texttt{Sufficiency}) & an objective-specific
    argument per model family \\
    Can the state be recovered? & $g_m(F_m(H))$ & readout bounds
    (\texttt{Measurement}, \texttt{Recovery}) & held-out calibration against
    the target's own reliability \\
    What changes financially? & $W$, $\Sigma$, scores & stability and
    likelihood (\texttt{Stability}, \texttt{Likelihood}) & matched
    econometric comparisons \\
    \bottomrule
  \end{tabular}
  \caption{Four questions a learned proximity measure must answer separately.}%
  \label{tab:framework}
\end{table}

\paragraph{Target.} Here $s_{\mathcal C}(H)$ is the state on which a declared
class $\mathcal C$ of economic responses depends,
for example the exposures $b(S_i)$ in the residual model above. The current
empirical stand-in is statistical: each firm's loadings on the leading principal
components of prior factor residuals of firms \emph{outside} the evaluated
cohort. It is a common-risk coordinate, not an identified economic exposure.
Its split-half reliability is reported with every decoder, because a noisy
target caps the attainable fit whatever the representation contains.

\paragraph{Recovery.} For a $K$-Lipschitz readout $g$, the measured state misses
$s$ by at most $K\,d(\widehat F,F)$ plus the alignment error
$\xi=d(g(F),s)$ (Theorem~\ref{thm:measurement-error}). That error belongs to
the pair $(F,g)$, not to $F$ alone.

\begin{proposition}[Information and readout split]\label{prop:readout-split}
  \lean{FreeFoundations.Recovery.excess_risk_condExp}\leanok%
  \uses{thm:recovery-excess-condexp}
  For a square-integrable state $s$ and a square-integrable readout $g(F)$,
  \begin{equation}\label{eq:readout-split}
    \E\norm{s-g(F)}^2
    =\underbrace{\E\norm{s-\E[s\mid
    F]}^2}_{\text{information absent from }F}
    +\underbrace{\E\norm{\E[s\mid
    F]-g(F)}^2}_{\text{readout misspecification or estimation}}.
  \end{equation}
\end{proposition}

The Lean statement is for a real state; a Euclidean state splits coordinate by
coordinate. A restricted decoder can therefore fail even when $F$
carries the state, and a
failed recovery experiment does not by itself show that the representation
lacks the information. The validation protocol fixes the target's coordinate
system before any fold is drawn, so a held-out firm's returns reach neither the
training targets nor the fitted readout. Fitting a readout to historical
returns for forecasting is legitimate, but it is a deployment procedure, not
this validation.

\paragraph{Financial use.} The expected Gaussian loss of the spatial-error
forecast is a variance term, a spectrum term and the filtered second moments
$m_i(\lambda)$ (Theorem~\ref{thm:likelihood-loss}).

\begin{proposition}[Neighbour reconstruction against imported
  noise]\label{prop:filtered}
  \lean{FreeFoundations.Likelihood.filtered_second_moment}\leanok%
  \uses{thm:likelihood-filtered}
  Under a common-plus-idiosyncratic residual model $u_i=a_i'q+e_i$, with
  $\operatorname{Var}(q)=I$ and independent $e_i$ of variance $\eta_i$, and a
  zero-diagonal network, the spatial filter's expected squared residual for
  firm $i$ is
  \begin{equation}
    m_i(\lambda)=\Big\lVert a_i-\lambda\sum_j W_{ij}a_j\Big\rVert^2+\eta_i
    +\lambda^2\sum_j W_{ij}^2\eta_j .
  \end{equation}
\end{proposition}

The network helps when a firm's neighbours reconstruct its common exposure, and
costs through the idiosyncratic noise it imports, weighted by the row's
concentration $\sum_j {W_{ij}}^2$. Neither a diffuse nor a
concentrated network is
favoured in general. The empirical split of the realized likelihood is exact
algebra, but it uses total residual second moments, which mix the common and
idiosyncratic parts, and it conditions on realized factors. A larger linear
(alignment) term therefore shows stronger weighted residual cross-products
under the fitted specification, not recovered exposures.

\subsection{One complete argument: TabPFN's late readout}
Each link is marked as published theory, a result formalized here, a
maintained assumption, or an empirical test.
\begin{enumerate}
  \item \emph{Objective to predictive statistic (published; formalized).}
    Prior-data fitted networks minimize the expected negative log-likelihood of
    held-out labels over datasets drawn from a synthetic prior. The excess over
    the Bayes loss is the expected Kullback-Leibler divergence to the posterior
    predictive under that prior~\cite{pfn}
    (Theorem~\ref{thm:score-pretraining}).
    The late query state feeds the output head, so the predictive distribution
    of the support label is a function of that state.
  \item \emph{Prior to firm data (assumed; probed).} The synthetic
    prior is not a
    model of disclosures. Read frequentistically, a PFN's variance vanishes as
    the in-context sample grows when its sensitivity to single support rows
    does, while its bias vanishes only if the predictor localizes around the
    query~\cite{nagler}. No mismatch rate is claimed. The support-size
    experiment varies the in-context sample on fixed query firms, which bears on
    the variance channel only.
  \item \emph{Predictive statistic to economic state (assumed; partly
    refuted).} Sufficiency requires the target state to be a function of the
    support label's conditional law (Theorem~\ref{thm:sufficiency-iff}). One
    scalar label cannot distinguish every multidimensional exposure
    (Theorem~\ref{thm:recovery-scalar}), so a nonzero information term
    in~\eqref{eq:readout-split} is expected.
  \item \emph{State to measured distance (formalized; tested).} The ridge
    readout is linear and Lipschitz, and its bound is
    Theorem~\ref{thm:measurement-error}. The test is held-out decoder fit
    against the target's reliability.
  \item \emph{Distance to covariance (formalized; conditional).}
    Proposition~\ref{prop:stability}, under its resolvent assumption.
\end{enumerate}
An architecture-matched control re-initializes the same checkpoint's weights
before encoding. Only for TabPFN does the comparison then isolate trained
weights; elsewhere it compares pretrained-model procedures.

\paragraph{Admissibility of other readouts.} TabFM and EXAONE enter as
comparators, not as instances of the argument above. Their training objectives
are taken from provider documentation and are not verified here. EXAONE's
task readout precedes a quantile output head: if that head is trained with a
pinball loss, the state determines predicted quantiles, a consistent target
for that loss~\cite{gneiting}, rather than the full conditional law. State
readouts taken before in-context prediction have no objective-specific
argument at all. No contrastive or predictive-latent objective is treated as a
proper scoring rule for an observed conditional law without a separate
argument.

\section{Data expansion and chronological design}
The acquisition target is 1,000 distinct current US exchange-listed SEC issuers,
Yahoo price histories from 2008 through 2025, and annual filing snapshots from
2009 through 2024. These are targets, not asserted usable sample sizes.
The first XBRL mandate dates to 2009~\cite{secapi}; realized coverage determines
the feasible evaluation years. A fact from an older fiscal period that first
appears in a later filing is not backfilled to an earlier information set.
Current-listing and capitalization selection introduces survivorship bias;
compressed source snapshots and exclusion logs make this explicit, but do not
remove it. Historical listings, delisting returns, and licensed accounting data
would be a separate data upgrade if available.

The main panel is unbalanced across years. Requiring complete ten-year histories
for every characteristic reduced the earlier 100-issuer exercise to six firms.
Strict debt-maturity mapping is particularly restrictive. Missing lease or debt
components must not be silently set to zero, and liabilities must not replace
debt. Report sample losses by year, characteristic, and mapping rule. Retain a
strict original-method cohort alongside the broader disclosure sample;
comparisons across different samples must not be presented as likelihood gains.

At each cutoff, determine eligibility using information already available.
Fit vocabulary selection, imputation, scaling, PCA, and autoencoders on past
observations; then encode the current disclosures and freeze $W_t$ before the
subsequent return window. Record foundation-model checkpoint dates, training
provenance where known, support labels, and preprocessing. Modern checkpoints
applied to historical data do not establish historical deployability.
Do not select firms using complete future return histories. Missing outcome
coordinates require a common observed sample and the corresponding marginal
model density, rather than silently rebuilding the network after
outcomes arrive.

\section{Single-parameter panel comparison}
For each encoder and input set, estimate
\begin{equation}
  r_t=B_t f_t+u_t,\qquad u_t=\lambda W_t u_t+\varepsilon_t,
  \qquad\varepsilon_t\sim N(0,D_t).
\end{equation}
Economic weights can be endogenous~\cite{qulee}. Chronological construction
prevents look-ahead but does not by itself identify causal spillovers.
We report two innovation specifications under the same estimation rule for
every network: the paper-aligned $D_t=\sigma_t^2I$, and a heterogeneous diagonal
$D_t$ with one profiled variance per firm and year. At $\lambda=0$ the latter
is exactly the diagonal factor benchmark, so differences from it isolate the
network. A third, variance-matched comparison imposes prior sample variances on
every forecast covariance, so that only implied correlations differ.
There is one scalar $\lambda$ per approach in a pooled estimation window;
$W_t$ can change annually. Fixed, predetermined refitting windows must be
identical across approaches. The main measurement comparison holds $B_t$,
$D_t$, factor definitions, dates, and firms fixed. CAPM aligns with
the reference;
FF5 plus momentum is a separate, explicitly labeled extension.

For raw sample-covariance comparisons, the dimension must be compatible
with the estimation window. With $N$ firms and $T$ demeaned daily observations,
the sample covariance has rank at most $T-1$. A 1,000-stock covariance from
approximately three years of returns is therefore singular. Prespecify a
longer window or a smaller common cohort for an unregularized benchmark;
otherwise label regularization explicitly rather than treating an undefined
Gaussian baseline density as a model victory.

Report training ML fit and aggregate subsequent held-out log-likelihood
separately. Conditional residual likelihood may use realized factors as
regressors; it is not an unconditional return forecast. For a forecast claim,
factor uncertainty must also be modeled using information available beforehand.
Use time blocks for uncertainty in paired likelihood differences, with
prespecified handling of overlapping estimation windows. Distinct networks
are generally nonnested; their likelihood difference is not automatically a
chi-squared likelihood-ratio test. Boundary estimates require explicit flags.

The ladder contains the individual reference networks, their predetermined
joint-distance control, direct and PCA/autoencoder controls, and pretrained
encoders. Each learned vector combines the entire selected feature set.
Compare four characteristic histories first, then the same histories plus
broader disclosures. Additional features and pretraining must be isolated in
separate comparisons. A fixed reference encoder/context check helps distinguish
changing firm states from changes in preprocessing or model context.

\section{Evidence needed before writing conclusions}
The first main table should document attainable coverage and missingness.
The second should report the matched likelihood comparison and uncertainty.
One figure can show the time pattern of paired log-score differences and
network concentration. Industry, size, leverage, reporting-pattern controls,
and fixed-reference measurement checks should address alternative explanations.

The measurement claim needs three further pieces of evidence, in this order.
First, held-out recovery: decoder fit against a target whose coordinates no
validation firm helped estimate, reported beside that target's split-half
reliability. Second, a reduced comparison on one input set. It should include
a direct-distance baseline, PCA, a supervised network trained on the same
label, a label-null version of the pretrained readout, and an untrained copy of
the same architecture where one can be built. Third, the same comparison at
several nested support sizes on fixed query firms, which tests whether any
pretrained advantage is a small-sample effect. Without an architecture-matched
control, the result is a comparison of pretrained-model procedures, not an
effect of pretraining. Economic validation against identified shock exposures
remains a separate step.
All already inspected pilot years are development evidence. Any subsequently
reserved evaluation block must be described honestly; this is not a claim of
preregistration or a previously untouched holdout.

The earlier six-firm fits and boundary estimates are diagnostic results. The
cross-sectional replication is now estimated on one common cohort per
characteristic set from the 1,000-issuer universe, with two-stage least squares
and maximum-likelihood standard errors, exact residual Moran moments,
conditional information criteria and a generated-beta bootstrap. Information
criteria rank networks with equal parameter counts exactly as their likelihoods
do, so they are reported for nonspatial versus spatial and single-factor versus
multifactor comparisons. None of these results supports a claim of superior
economic measurement, nor can a negative spatial coefficient be relabeled
positive shock propagation.
If gains do not survive matched-input and chronological controls, the paper
must report that finding or revise its question rather than select a favorable
model after examining evaluation returns.

\appendix
\section{Proof of the conditional stability result}
The resolvent identity gives
$A-A^\star=\lambda A(W-W^\star)A^\star$, hence
$\norm{A-A^\star}_2\leq |\lambda|M^2\norm{W-W^\star}_2$.
Expanding the covariance difference as
$(A-A^\star)DA'+A^\star D(A-A^\star)'$ and using submultiplicativity
establishes the first bound.
For the second, the ratio of corresponding positive kernel entries lies
between $e^{-\epsilon/h}$ and $e^{\epsilon/h}$; the same holds for their
row-sum ratio. Normalized entry ratios therefore lie between
$e^{-2\epsilon/h}$ and $e^{2\epsilon/h}$. Summing absolute entry differences
against a probability row bounds the row sum by $e^{2\epsilon/h}-1$;
the distance between any two probability rows is also at most two.

% Formal results of the Lean library, one statement per declaration.
% Shared by the web blueprint, the printed blueprint and the manuscript appendix.
% Every \lean{...} name is checked against the compiled library by
% `just blueprint-check`; \leanok marks a statement or proof with no `sorry`.

\section{Reading this blueprint}\label{sec:bp-reading}

Each statement below restates one declaration of the \texttt{FreeFoundations}
library in ordinary notation. Each proof is a short account of the formal
argument. The results follow the library's four questions: what the economic
target is (\S\ref{sec:bp-target}), what a training objective makes an optimal
model keep (\S\ref{sec:bp-score}, \S\ref{sec:bp-alignment}), whether the state
and its distances can be read out (\S\ref{sec:bp-recovery},
\S\ref{sec:bp-geometry}), and what the spatial model does with the distances
(\S\ref{sec:bp-network}, \S\ref{sec:bp-likelihood}, \S\ref{sec:bp-stability}).
Each result holds under its stated hypotheses only. None of them shows that a
particular checkpoint satisfies those hypotheses.

% ---------------------------------------------------------------------------
\section{The economic target}\label{sec:bp-target}

Residual returns contain common components $b_i'q$ with shock second moment
$\Omega$. In whitened coordinates $a_i=\Omega^{1/2}b_i$.

\begin{definition}[Quadratic form]\label{def:exposure-form}
  \lean{FreeFoundations.Exposure.form}\leanok
  $\operatorname{form}_\Omega(x,y)=x'\Omega y$.
\end{definition}

\begin{lemma}[Second moment]\label{thm:exposure-covariance}
  \lean{FreeFoundations.Exposure.covariance_eq}\leanok
  \uses{def:exposure-form}
  For symmetric $\Omega$,
  \[
    b_i'\Omega b_j=\tfrac12\bigl(b_i'\Omega b_i+b_j'\Omega b_j
    -(b_i-b_j)'\Omega(b_i-b_j)\bigr).
  \]
\end{lemma}
\begin{proof}\leanok
  Expand the quadratic form of $b_i-b_j$ and use $\Omega'=\Omega$.
\end{proof}

\begin{lemma}[Whitened second moment]\label{thm:exposure-inner}
  \lean{FreeFoundations.Exposure.inner_eq_half}\leanok
  In any real inner product space,
  $\ip{a_i}{a_j}=\tfrac12(\norm{a_i}^2+\norm{a_j}^2-\norm{a_i-a_j}^2)$.
  Distance alone does not order covariance: the common variances matter too.
\end{lemma}
\begin{proof}\leanok
  The polarization identity.
\end{proof}

\begin{lemma}[Pricing errors are Lipschitz in exposure]\label{thm:exposure-pricing}
  \lean{FreeFoundations.Exposure.pricing_error_lipschitz}\leanok
  $\abs{\ip{a_i}{g}-\ip{a_j}{g}}\le\norm{a_i-a_j}\,\norm{g}$.
  If the omitted factors are traded and firms earn no other expected return,
  close exposures imply close factor-implied expected returns, with
  $\norm{g}=(\gamma'\Omega^{-1}\gamma)^{1/2}$ the maximum Sharpe ratio of the
  factor prices $\gamma$.
\end{lemma}
\begin{proof}\leanok
  Cauchy--Schwarz applied to $\ip{a_i-a_j}{g}$.
\end{proof}

\begin{lemma}[Exposure distance is a response variance]\label{thm:score-response}
  \lean{FreeFoundations.Score.response_sq}\leanok
  \uses{def:exposure-form, def:score-second-moment}
  For shocks $q$ with finite second moment $\Omega=\E[qq']$ and any $c$,
  $\E[(c'q)^2]=c'\Omega c$. With $c=b_i-b_j$, the squared exposure distance is
  the second moment of the two firms' response difference.
\end{lemma}
\begin{proof}\leanok
  Expand $(c'q)^2$ as a double sum and integrate term by term.
\end{proof}

% ---------------------------------------------------------------------------
\section{What scores measure}\label{sec:bp-score}

\begin{lemma}[Quadratic forms are traces]\label{thm:score-trace}
  \lean{FreeFoundations.Score.quadratic_eq_trace}\leanok
  $u'Au=\tr(A\,uu')$.
\end{lemma}
\begin{proof}\leanok
  Expand both sides entrywise.
\end{proof}

\begin{lemma}[Realized score]\label{thm:score-realized}
  \lean{FreeFoundations.Score.realized_quadratic}\leanok
  \uses{thm:score-trace}
  $\sum_t u_t'Au_t=\tr\bigl(A\sum_t u_tu_t'\bigr)$: the summed quadratic terms
  depend on the observations only through their summed outer products.
\end{lemma}
\begin{proof}\leanok
  Sum Lemma~\ref{thm:score-trace} and use linearity of the trace.
\end{proof}

\begin{definition}[Second-moment matrix]\label{def:score-second-moment}
  \lean{FreeFoundations.Score.secondMoment}\leanok
  $\E[uu']$, entrywise $\int u_iu_j\,d\mu$.
\end{definition}

\begin{lemma}[Expected quadratic form]\label{thm:score-integral}
  \lean{FreeFoundations.Score.integral_quadratic}\leanok
  \uses{def:score-second-moment}
  With finite second moments, $\E[u'Au]=\tr(A\,\E[uu'])$.
\end{lemma}
\begin{proof}\leanok
  Integrate the finite double sum term by term.
\end{proof}

\begin{definition}[Gaussian and Stein losses]\label{def:score-loss}
  \lean{FreeFoundations.Score.loss, FreeFoundations.Score.stein}\leanok
  $L(C,S)=\log\det C+\tr(C^{-1}S)$ and
  $\operatorname{Stein}(C,S)=\tr(C^{-1}S)-\log\det S+\log\det C-N$.
\end{definition}

\begin{lemma}[Expected Gaussian loss]\label{thm:score-expected-loss}
  \lean{FreeFoundations.Score.expected_loss}\leanok
  \uses{def:score-loss, def:score-second-moment}
  With finite second moments $M$, a covariance forecast $C$ has expected loss
  \[\log\det C+\E[u'C^{-1}u]=L(C,M).\]
  The law of $u$ enters only through $M$;
  no normality is assumed.
\end{lemma}
\begin{proof}\leanok
  \uses{thm:score-integral}
  Lemma~\ref{thm:score-integral} with $A=C^{-1}$.
\end{proof}

\begin{lemma}[Stein's loss is nonnegative]\label{thm:score-stein}
  \lean{FreeFoundations.Score.stein_nonneg}\leanok
  \uses{def:score-loss}
  For positive definite $C$ and $S$, $\operatorname{Stein}(C,S)\ge0$.
\end{lemma}
\begin{proof}\leanok
  Factor $C=B'B$. Then $M=B'^{-1}SB^{-1}$ is positive definite with
  $\tr M=\tr(C^{-1}S)$ and $\det M=\det S/\det C$, so Stein's loss is
  $\sum_k(\mu_k-\log\mu_k-1)\ge0$ over the eigenvalues $\mu_k>0$ of $M$.
\end{proof}

\begin{lemma}[Loss gap]\label{thm:score-loss-sub}
  \lean{FreeFoundations.Score.loss_sub}\leanok
  \uses{def:score-loss}
  $L(C,S)-L(S,S)=\operatorname{Stein}(C,S)$ for positive definite $S$.
\end{lemma}
\begin{proof}\leanok
  $\tr(S^{-1}S)=N$ and $\log\det C-\log\det S=-\log\det(C^{-1}S)$.
\end{proof}

\begin{theorem}[The Gaussian score is proper for the second moment]\label{thm:score-proper}
  \lean{FreeFoundations.Score.loss_le}\leanok
  \uses{def:score-loss}
  $L(S,S)\le L(C,S)$ for positive definite $C$ and $S$.
\end{theorem}
\begin{proof}\leanok
  \uses{thm:score-stein, thm:score-loss-sub}
  The gap is Stein's loss, which is nonnegative.
\end{proof}

\begin{definition}[Cross-entropy and KL divergence]\label{def:score-entropy}
  \lean{FreeFoundations.Score.crossEntropy, FreeFoundations.Score.kl}\leanok
  On a finite outcome space, $H(p,q)=-\sum_y p_y\log q_y$ and
  $\KL(p\Vert q)=\sum_y p_y\log(p_y/q_y)$.
\end{definition}

\begin{lemma}[Cross-entropy is entropy plus KL]\label{thm:score-ce-eq}
  \lean{FreeFoundations.Score.crossEntropy_eq}\leanok
  \uses{def:score-entropy}
  For $q>0$, $H(p,q)=H(p,p)+\KL(p\Vert q)$.
\end{lemma}
\begin{proof}\leanok
  Split $\log q_y=\log p_y-\log(p_y/q_y)$ where $p_y>0$; terms with $p_y=0$
  vanish on both sides.
\end{proof}

\begin{lemma}[Gibbs' inequality]\label{thm:score-kl}
  \lean{FreeFoundations.Score.kl_nonneg}\leanok
  \uses{def:score-entropy}
  $\KL(p\Vert q)\ge0$ for $p\ge0$, $q>0$ with equal total mass.
\end{lemma}
\begin{proof}\leanok
  $\log x\le x-1$ applied to $q_y/p_y$.
\end{proof}

\begin{theorem}[The log score is proper]\label{thm:score-log-proper}
  \lean{FreeFoundations.Score.crossEntropy_le}\leanok
  \uses{def:score-entropy}
  $H(p,p)\le H(p,q)$: the true law has the smallest expected log loss.
\end{theorem}
\begin{proof}\leanok
  \uses{thm:score-ce-eq, thm:score-kl}
  Lemmas~\ref{thm:score-ce-eq} and~\ref{thm:score-kl}.
\end{proof}

\begin{theorem}[Pretraining excess loss]\label{thm:score-pretraining}
  \lean{FreeFoundations.Score.pretraining_excess}\leanok
  \uses{def:score-entropy}
  Averaged over contexts $c$ with weights $\pi_c$, the log loss of a predictor
  $q$ exceeds the Bayes loss of the posterior predictive $p$ by the expected KL:
  \[
    \sum_c\pi_cH(p_c,q_c)-\sum_c\pi_cH(p_c,p_c)=\sum_c\pi_c\KL(p_c\Vert q_c).
  \]
  For a prior-data fitted network, $\pi$ is the synthetic task prior and $p_c$
  the posterior predictive given support set $c$.
\end{theorem}
\begin{proof}\leanok
  \uses{thm:score-ce-eq}
  Apply Lemma~\ref{thm:score-ce-eq} in each context and sum.
\end{proof}

\begin{lemma}[Empirical risk minimisation]\label{thm:score-erm}
  \lean{FreeFoundations.Score.erm_bound}\leanok
  If $\abs{\widehat R(\psi)-R(\psi)}\le g$ for every $\psi$ and $\widehat\psi$
  minimises $\widehat R$ to within $\varepsilon$, then
  $R(\widehat\psi)\le R(\psi)+2g+\varepsilon$ for every $\psi$.
\end{lemma}
\begin{proof}\leanok
  Chain the three inequalities.
\end{proof}

% ---------------------------------------------------------------------------
\section{Alignment and sufficiency}\label{sec:bp-alignment}

\begin{definition}[Optimal and sufficient representations]\label{def:alignment}
  \lean{FreeFoundations.Alignment.Optimal, FreeFoundations.Alignment.Sufficient}\leanok
  A representation $Z$ of histories $H$ is \emph{optimal} for a target
  prediction $F$ if $F$ factors through $Z$, and \emph{sufficient} for the
  economic state if the state factors through $Z$.
\end{definition}

\begin{theorem}[Deterministic alignment]\label{thm:alignment-iff}
  \lean{FreeFoundations.Alignment.optimal_imp_sufficient_iff}\leanok
  \uses{def:alignment}
  Every target-optimal representation is sufficient if and only if
  $F(h)=F(h')$ implies $\mathrm{state}(h)=\mathrm{state}(h')$: the target
  separates every pair of histories the economic state separates.
\end{theorem}
\begin{proof}\leanok
  For the forward direction take $Z=F$, which is optimal. Conversely an optimal
  $Z$ determines $F$, which determines the state.
\end{proof}

\begin{lemma}[When the state itself is optimal]\label{thm:alignment-state}
  \lean{FreeFoundations.Alignment.state_optimal_iff}\leanok
  \uses{def:alignment}
  $F=d\circ\mathrm{state}$ for some $d$ if and only if the target needs no
  distinction the state ignores.
\end{lemma}
\begin{proof}\leanok
  Choose $d$ on the range of the state; it is well defined by the hypothesis.
\end{proof}

\begin{lemma}[Label-null targets lose the state]\label{thm:alignment-null}
  \lean{FreeFoundations.Alignment.constant_target_loses_state}\leanok
  \uses{def:alignment}
  With a constant target, the representation that keeps nothing is optimal,
  and it is not sufficient whenever the state varies.
\end{lemma}
\begin{proof}\leanok
  A constant factors through the trivial representation; two histories with
  different states are then merged.
\end{proof}

\begin{lemma}[Marginals need not carry the joint law]\label{thm:alignment-marginals}
  \lean{FreeFoundations.Alignment.pair, FreeFoundations.Alignment.marginals_free_of_state, FreeFoundations.Alignment.joint_depends_on_state}\leanok
  Let $Y_1=R$ and $Y_2=R\oplus S$ for a fair coin $R$ and an observed bit $S$.
  Both marginals are fair coins whatever $S$, yet
  $P(Y_1=Y_2=1)$ is $1/2$ when $S=0$ and $0$ when $S=1$.
  Exact marginal forecasts can discard the state that fixes dependence.
\end{lemma}
\begin{proof}\leanok
  Direct computation on the four outcomes.
\end{proof}

\begin{lemma}[Attaining the Bayes predictor]\label{thm:sufficiency-condexp}
  \lean{FreeFoundations.Sufficiency.condExp_eq_iff}\leanok
  For $\sigma$-algebras $\mathcal Z\subseteq\mathcal H$,
  $\E[X\mid\mathcal Z]=\E[X\mid\mathcal H]$ almost surely if and only if
  $\E[X\mid\mathcal H]$ is almost surely $\mathcal Z$-measurable.
\end{lemma}
\begin{proof}\leanok
  The tower property in one direction; in the other, a $\mathcal Z$-measurable
  function is its own conditional expectation given $\mathcal Z$.
\end{proof}

\begin{theorem}[Statistical alignment]\label{thm:sufficiency-iff}
  \lean{FreeFoundations.Sufficiency.sufficient_iff}\leanok
  Every representation that reproduces the pretraining Bayes predictor
  $\E[X\mid H]$ also reproduces the economic one $\E[Y\mid H]$ if and only if
  $\E[Y\mid H]=g(\E[X\mid H])$ almost surely for a measurable $g$. Conditional
  laws follow from indicator targets.
\end{theorem}
\begin{proof}\leanok
  \uses{thm:sufficiency-condexp}
  Take $\mathcal Z=\sigma(\E[X\mid H])$ and apply
  Lemma~\ref{thm:sufficiency-condexp}; measurability with respect to
  $\sigma(\E[X\mid H])$ is the Doob--Dynkin factorization.
\end{proof}

\begin{corollary}[Label-null, statistically]\label{thm:sufficiency-constant}
  \lean{FreeFoundations.Sufficiency.constant_target_iff}\leanok
  If the pretraining prediction is almost surely constant, the objective
  guarantees economic sufficiency only for an almost surely constant economic
  prediction.
\end{corollary}
\begin{proof}\leanok
  \uses{thm:sufficiency-iff}
  A function of a constant is constant.
\end{proof}

% ---------------------------------------------------------------------------
\section{Recovery and measurement}\label{sec:bp-recovery}

Tasks are a linear map $C$ from exposures to task means; a decoder $L$ reads
exposures back.

\begin{definition}[Distinguishing tasks and decoders]\label{def:recovery}
  \lean{FreeFoundations.Recovery.Distinguishes, FreeFoundations.Recovery.Decoder}\leanok
  $C$ distinguishes exposures with margin $\kappa$ if
  $\norm{Ca}\ge\kappa\norm{a}$ for every $a$. A decoder is a left inverse $L$
  of $C$ with $\kappa\norm{Lt}\le\norm{t}$ for every $t$.
\end{definition}

\begin{lemma}[Distinguishing tasks are injective]\label{thm:recovery-injective}
  \lean{FreeFoundations.Recovery.Distinguishes.injective}\leanok
  \uses{def:recovery}
  If $\kappa>0$, a distinguishing $C$ is injective.
\end{lemma}
\begin{proof}\leanok
  $Ca=0$ forces $\kappa\norm a\le0$.
\end{proof}

\begin{lemma}[Dimension]\label{thm:recovery-finrank}
  \lean{FreeFoundations.Recovery.finrank_le_of_distinguishes}\leanok
  \uses{def:recovery}
  Distinguishing tasks need at least as many dimensions as the exposures.
\end{lemma}
\begin{proof}\leanok
  \uses{thm:recovery-injective}
  An injective linear map does not lower dimension.
\end{proof}

\begin{theorem}[One target is not enough]\label{thm:recovery-scalar}
  \lean{FreeFoundations.Recovery.scalar_not_distinguishes}\leanok
  \uses{def:recovery}
  A single scalar task cannot distinguish exposures of dimension two or more,
  however accurately it is predicted.
\end{theorem}
\begin{proof}\leanok
  \uses{thm:recovery-finrank}
  The target space has dimension one.
\end{proof}

\begin{lemma}[An aggregate target merges firms]\label{thm:recovery-aggregate}
  \lean{FreeFoundations.Recovery.Plane, FreeFoundations.Recovery.aggregate_merges}\leanok
  The target $a_f+a_o$ takes the same value on $(1,0)$ and $(0,1)$, which are
  $\sqrt2$ apart.
\end{lemma}
\begin{proof}\leanok
  Direct computation.
\end{proof}

\begin{lemma}[Diverse targets distinguish]\label{thm:recovery-diverse}
  \lean{FreeFoundations.Recovery.diverse_distinguishes}\leanok
  $(a_0+a_1)^2+(a_0-a_1)^2=2\norm a^2$, so the two targets $a_f\pm a_o$
  distinguish exposures with margin $\sqrt2$.
\end{lemma}
\begin{proof}\leanok
  Expand the squares.
\end{proof}

\begin{theorem}[A decoder exists]\label{thm:recovery-decoder}
  \lean{FreeFoundations.Recovery.exists_decoder}\leanok
  \uses{def:recovery}
  In finite dimension, tasks that distinguish exposures with margin
  $\kappa>0$ admit a decoder with gain at most $1/\kappa$.
\end{theorem}
\begin{proof}\leanok
  \uses{thm:recovery-injective}
  Project orthogonally onto the range of $C$, then invert $C$ there.
\end{proof}

\begin{lemma}[Exposure error]\label{thm:recovery-exposure}
  \lean{FreeFoundations.Recovery.exposure_error}\leanok
  \uses{def:recovery}
  $\kappa\norm{L\hat\tau-a}\le\norm{\hat\tau-Ca}$.
\end{lemma}
\begin{proof}\leanok
  $L\hat\tau-a=L(\hat\tau-Ca)$ and $L$ has gain $1/\kappa$.
\end{proof}

\begin{lemma}[Exposure error with mismatch]\label{thm:recovery-mismatch}
  \lean{FreeFoundations.Recovery.exposure_error_mismatch, FreeFoundations.Recovery.exposure_error_mismatch_sq}\leanok
  \uses{def:recovery}
  If the conditional task mean is $\tau=Ca+v$, then
  $\kappa\norm{L\hat\tau-a}\le\norm{\hat\tau-\tau}+\norm v$ and
  $\kappa^2\norm{L\hat\tau-a}^2\le2(\norm{\hat\tau-\tau}^2+\norm v^2)$.
  Forecast accuracy cannot remove $v$.
\end{lemma}
\begin{proof}\leanok
  \uses{thm:recovery-exposure}
  The triangle inequality, then $(x+y)^2\le2x^2+2y^2$.
\end{proof}

\begin{lemma}[In mean square]\label{thm:recovery-mean}
  \lean{FreeFoundations.Recovery.mean_exposure_error}\leanok
  \uses{def:recovery}
  For any measure over firms or information sets,
  $\kappa^2\E\norm{L\hat\tau-a}^2\le\E\norm{\hat\tau-Ca}^2$. A finite sample is
  the counting measure.
\end{lemma}
\begin{proof}\leanok
  \uses{thm:recovery-exposure}
  Square Lemma~\ref{thm:recovery-exposure} and integrate.
\end{proof}

\begin{lemma}[Distance error]\label{thm:recovery-distance}
  \lean{FreeFoundations.Recovery.distance_error}\leanok
  $\bigl\lvert\norm{\hat a_i-\hat a_j}-\norm{a_i-a_j}\bigr\rvert
  \le\norm{\hat a_i-a_i}+\norm{\hat a_j-a_j}$.
\end{lemma}
\begin{proof}\leanok
  The reverse triangle inequality.
\end{proof}

\begin{theorem}[Distance error over all pairs]\label{thm:recovery-pairwise}
  \lean{FreeFoundations.Recovery.pairwise_distance_error}\leanok
  Summed over ordered pairs of distinct firms,
  \[
    \sum_i\sum_{j\ne i}(\hat d_{ij}-d_{ij})^2\le4(N-1)\sum_i\norm{\hat a_i-a_i}^2.
  \]
  The inequality is deterministic, so dependence across pairs is irrelevant.
\end{theorem}
\begin{proof}\leanok
  \uses{thm:recovery-distance}
  Square Lemma~\ref{thm:recovery-distance}, bound $(x+y)^2\le2x^2+2y^2$, and
  count how often each firm appears.
\end{proof}

\begin{lemma}[Excess risk is a projection distance]\label{thm:recovery-excess}
  \lean{FreeFoundations.Recovery.excess_risk}\leanok
  If $\tau$ is the orthogonal projection of $Y$ onto a closed subspace $K$ and
  $f\in K$, then $\norm{f-Y}^2-\norm{\tau-Y}^2=\norm{f-\tau}^2$.
\end{lemma}
\begin{proof}\leanok
  Pythagoras: $f-\tau\in K$ is orthogonal to $Y-\tau$.
\end{proof}

\begin{theorem}[Information and readout split]\label{thm:recovery-excess-condexp}
  \lean{FreeFoundations.Recovery.excess_risk_condExp}\leanok
  For a square-integrable real $Y$ and any $\mathcal F$-measurable predictor $f$
  in $L^2$,
  \[
    \E(Y-f)^2=\E\bigl(Y-\E[Y\mid\mathcal F]\bigr)^2
    +\E\bigl(\E[Y\mid\mathcal F]-f\bigr)^2 .
  \]
  With $Y=s$ the target state, $\mathcal F=\sigma(F)$ the representation and
  $f=g(F)$ a readout, the first term is information that $F$ does not carry and
  the second is readout misspecification or estimation error. A Euclidean state
  splits coordinate by coordinate.
\end{theorem}
\begin{proof}\leanok
  \uses{thm:recovery-excess}
  Lemma~\ref{thm:recovery-excess} with $K$ the $\mathcal F$-measurable
  subspace of $L^2$, onto which conditional expectation is the orthogonal
  projection.
\end{proof}

\begin{lemma}[Readouts are reparametrization invariant]\label{thm:recovery-invariant}
  \lean{FreeFoundations.Recovery.readout_invariant}\leanok
  If a linear readout $G$ gives exposures from states $z$, the reparametrized
  states $ez$ with readout $Ge^{-1}$ give the same exposures and distances.
  Native distances between states have no such invariance
  (\S\ref{sec:bp-geometry}).
\end{lemma}
\begin{proof}\leanok
  $Ge^{-1}(ez)=Gz$.
\end{proof}

\begin{theorem}[Measurement error]\label{thm:measurement-error}
  \lean{FreeFoundations.Measurement.measurement_error}\leanok
  For a $K$-Lipschitz readout $g$ from predictive statistics to states,
  \[
    d\bigl(g(\hat F),s\bigr)\le K\,d(\hat F,F)+\xi,\qquad\xi=d\bigl(g(F),s\bigr).
  \]
  The alignment error $\xi$ belongs to the pair $(F,g)$; Theorem
  \ref{thm:recovery-excess-condexp} splits its mean square.
\end{theorem}
\begin{proof}\leanok
  The triangle inequality through $g(F)$, then the Lipschitz bound.
\end{proof}

\begin{corollary}[In root mean square]\label{thm:measurement-rms}
  \lean{FreeFoundations.Measurement.rms_measurement_error}\leanok
  Under any measure,
  $\norm{d(g(\hat F),s)}_2\le K\norm{d(\hat F,F)}_2+\norm{\xi}_2$.
\end{corollary}
\begin{proof}\leanok
  \uses{thm:measurement-error}
  Theorem~\ref{thm:measurement-error} pointwise, then Minkowski's inequality in
  $L^2$.
\end{proof}

\begin{lemma}[Linear decoders are Lipschitz]\label{thm:measurement-decoder}
  \lean{FreeFoundations.Measurement.decoder_lipschitz}\leanok
  \uses{def:recovery}
  A decoder with margin $\kappa>0$ is $\kappa^{-1}$-Lipschitz, so the
  recovery bounds are Theorem~\ref{thm:measurement-error} with $K=1/\kappa$.
\end{lemma}
\begin{proof}\leanok
  Linearity and the gain bound.
\end{proof}

\begin{lemma}[Merged states]\label{thm:measurement-merged}
  \lean{FreeFoundations.Measurement.merged_error}\leanok
  Any single estimate $r$ satisfies $d(s,s')\le2\max(d(r,s),d(r,s'))$.
\end{lemma}
\begin{proof}\leanok
  The triangle inequality through $r$.
\end{proof}

\begin{theorem}[Lower bound]\label{thm:measurement-lower}
  \lean{FreeFoundations.Measurement.readout_error_ge}\leanok
  If a representation gives two histories the same value, every readout
  misses one of their states by at least half the distance between them,
  however accurate the prediction.
\end{theorem}
\begin{proof}\leanok
  \uses{thm:measurement-merged}
  Both readouts are the same point; apply Lemma~\ref{thm:measurement-merged}.
\end{proof}

\begin{corollary}[Reading two exposures from their aggregate]\label{thm:measurement-aggregate}
  \lean{FreeFoundations.Measurement.aggregate_readout_error}\leanok
  Any readout of the aggregate $a_f+a_o=1$ misses $(1,0)$ or $(0,1)$ by at
  least $\sqrt2/2$.
\end{corollary}
\begin{proof}\leanok
  \uses{thm:measurement-merged, thm:recovery-aggregate}
  The two exposures are $\sqrt2$ apart.
\end{proof}

\begin{theorem}[Posterior means understate distances]\label{thm:measurement-mean-sq}
  \lean{FreeFoundations.Measurement.mean_sq_distance, FreeFoundations.Measurement.sq_norm_mean_le}\leanok
  For a random difference $D=A_i-A_j$ with finite second moment,
  $\E\norm D^2=\norm{\E D}^2+\E\norm{D-\E D}^2$, so
  $\norm{\E D}^2\le\E\norm D^2$. Under the conditional law, distances between
  posterior-mean exposures understate expected distances by
  $\tr\operatorname{Var}(A_i-A_j\mid H)$, which differs across pairs.
\end{theorem}
\begin{proof}\leanok
  Expand $\norm{D}^2=\norm{\E D}^2+2\ip{\E D}{D-\E D}+\norm{D-\E D}^2$; the
  cross term integrates to zero.
\end{proof}

% ---------------------------------------------------------------------------
\section{Information is not geometry}\label{sec:bp-geometry}

\begin{lemma}[Reparametrization keeps optimality and sufficiency]\label{thm:geometry-equiv}
  \lean{FreeFoundations.Geometry.optimal_equiv, FreeFoundations.Geometry.sufficient_equiv}\leanok
  \uses{def:alignment}
  Any bijection $e$ of the representation space preserves both properties.
\end{lemma}
\begin{proof}\leanok
  Factor through $e^{-1}$.
\end{proof}

\begin{theorem}[Same information, different distances]\label{thm:geometry-rescale}
  \lean{FreeFoundations.Geometry.rescale}\leanok
  \uses{def:alignment}
  Scaling a representation by $c\ne0$ keeps it sufficient and multiplies every
  distance by $\abs c$.
\end{theorem}
\begin{proof}\leanok
  \uses{thm:geometry-equiv}
  Scaling is a bijection; $\norm{cz-cz'}=\abs c\norm{z-z'}$.
\end{proof}

\begin{theorem}[Observational equivalence]\label{thm:geometry-equivalence}
  \lean{FreeFoundations.Geometry.observationally_equivalent}\leanok
  Task means identify exposures only up to an invertible map:
  $(CG^{-1})(Ga)=Ca$.
\end{theorem}
\begin{proof}\leanok
  Cancel $G^{-1}G$.
\end{proof}

\begin{theorem}[Which maps preserve distances]\label{thm:geometry-distances}
  \lean{FreeFoundations.Geometry.distances_preserved_iff}\leanok
  A linear $G$ preserves every distance if and only if it preserves norms. The
  common covariance fixes the Gram matrix (Lemma~\ref{thm:exposure-inner}),
  which pins distances down up to rotation.
\end{theorem}
\begin{proof}\leanok
  Take $y=0$ in one direction; use $Gx-Gy=G(x-y)$ in the other.
\end{proof}

\begin{definition}[Chord baseline]\label{def:geometry-chord}
  \lean{FreeFoundations.Geometry.chord}\leanok
  $v(x)=[x,1]/\norm{[x,1]}$, the direct-distance baseline.
\end{definition}

\begin{lemma}[The chord map is injective]\label{thm:geometry-chord-injective}
  \lean{FreeFoundations.Geometry.chord_left_inverse, FreeFoundations.Geometry.chord_injective, FreeFoundations.Geometry.one_add_sq_pos, FreeFoundations.Geometry.chord_scale_pos}\leanok
  \uses{def:geometry-chord}
  $x$ is recovered as $v_{1:p}/v_{p+1}$, since $1+\sum_ix_i^2>0$.
\end{lemma}
\begin{proof}\leanok
  The common normalizing factor cancels in the ratio.
\end{proof}

\begin{lemma}[The chord map is bounded]\label{thm:geometry-chord-bounded}
  \lean{FreeFoundations.Geometry.chord_sq_sum, FreeFoundations.Geometry.chord_dist_sq_le}\leanok
  \uses{def:geometry-chord}
  $v(x)$ lies on the unit sphere, so every squared chord distance is at most 4.
  Injective yet bounded: the baseline keeps the information but compresses the
  geometry.
\end{lemma}
\begin{proof}\leanok
  $\norm{v-w}^2\le(\norm v+\norm w)^2=4$.
\end{proof}

% ---------------------------------------------------------------------------
\section{Kernels, networks and value}\label{sec:bp-network}

\begin{lemma}[Normalized ratios]\label{thm:kernel-normalized}
  \lean{FreeFoundations.Kernel.normalized_l1_le, FreeFoundations.Kernel.normalized_l1_le_two}\leanok
  Positive weights whose ratios lie in $[1/r,r]$ normalize to probability
  vectors within $\ell_1$ distance $r^2-1$. Any two normalizations are within
  $\ell_1$ distance 2.
\end{lemma}
\begin{proof}\leanok
  Each normalized ratio lies in $[r^{-2},r^2]$; sum the deviations against a
  probability vector. The second bound is the triangle inequality.
\end{proof}

\begin{definition}[Exponential kernel]\label{def:kernel-exp}
  \lean{FreeFoundations.Kernel.expRow, FreeFoundations.Kernel.expKernel}\leanok
  $W_{ij}=\exp(-d_{ij}/h)/\sum_{k\ne i}\exp(-d_{ik}/h)$ for $i\ne j$, and
  $W_{ii}=0$.
\end{definition}

\begin{lemma}[Scale]\label{thm:kernel-scale}
  \lean{FreeFoundations.Kernel.expRow_scale}\leanok
  \uses{def:kernel-exp}
  $(d,h)$ and $(cd,ch)$ give the same network for $c\ne0$.
\end{lemma}
\begin{proof}\leanok
  $cd/(ch)=d/h$.
\end{proof}

\begin{theorem}[Kernel sensitivity]\label{thm:kernel-sensitivity}
  \lean{FreeFoundations.Kernel.expRow_l1_le, FreeFoundations.Kernel.expKernel_linfty_le}\leanok
  \uses{def:kernel-exp}
  With at least two firms, if every off-diagonal distance moves by at most
  $\delta$ then $\norm{\widehat W-W}_\infty\le\min\{2,e^{2\delta/h}-1\}$.
\end{theorem}
\begin{proof}\leanok
  \uses{thm:kernel-normalized}
  Kernel ratios lie in $[e^{-\delta/h},e^{\delta/h}]$; apply
  Lemma~\ref{thm:kernel-normalized} to each row with $r=e^{\delta/h}$.
\end{proof}

\begin{lemma}[Exact expansion]\label{thm:network-expansion}
  \lean{FreeFoundations.Network.inverse_expansion}\leanok
  $(I-\lambda W)^{-1}=I+\lambda W+\lambda^2W^2(I-\lambda W)^{-1}$.
\end{lemma}
\begin{proof}\leanok
  Multiply through by $I-\lambda W$.
\end{proof}

\begin{lemma}[Inverse on an eigenvector]\label{thm:network-eigen}
  \lean{FreeFoundations.Network.inverse_on_eigenvector}\leanok
  If $(I-\lambda W)v=\mu v$ with $\mu\ne0$, then $(I-\lambda W)^{-1}v=\mu^{-1}v$.
\end{lemma}
\begin{proof}\leanok
  Apply the inverse to both sides.
\end{proof}

\begin{lemma}[Row-stochastic networks carry a common component]\label{thm:network-const}
  \lean{FreeFoundations.Network.rowStochastic_const, FreeFoundations.Network.rowStochastic_inverse_const}\leanok
  If $W\mathbf1=\mathbf1$ then $(I-\lambda W)^{-1}\mathbf1=\mathbf1/(1-\lambda)$
  for $\lambda\ne1$: the spatial covariance always contains a common component
  whose scale grows without bound as $\lambda\to1$.
\end{lemma}
\begin{proof}\leanok
  \uses{thm:network-eigen}
  $\mathbf1$ is an eigenvector of $I-\lambda W$ with eigenvalue $1-\lambda$.
\end{proof}

\begin{lemma}[The uniform network]\label{thm:network-uniform}
  \lean{FreeFoundations.Network.uniform, FreeFoundations.Network.uniform_const, FreeFoundations.Network.uniform_sum_zero}\leanok
  The uniform network $U$ with weight $1/(N-1)$ on every other firm fixes
  $\mathbf1$ and maps every vector summing to zero to $-1/(N-1)$ times itself:
  it is one equicorrelated factor.
\end{lemma}
\begin{proof}\leanok
  $Uv=(\mathbf1'v\,\mathbf1-v)/(N-1)$.
\end{proof}

\begin{lemma}[Placebos keep the spectrum and group mass]\label{thm:network-placebo}
  \lean{FreeFoundations.Network.placebo_det, FreeFoundations.Network.placebo_group_mass}\leanok
  Relabelling firms by a permutation leaves $\det(I-\lambda W)$, and so every
  likelihood's Jacobian term, unchanged. A within-group permutation preserves
  the total weight linking any two groups.
\end{lemma}
\begin{proof}\leanok
  Conjugation by a permutation matrix preserves the determinant; reindex the
  sums.
\end{proof}

\begin{theorem}[Hub plus proximity]\label{thm:value-decomposition}
  \lean{FreeFoundations.Value.decomposition}\leanok
  \uses{thm:exposure-inner}
  For row-stochastic $W$, doubly stochastic $U$ and common covariance
  $C_{ij}=\tfrac12(k_i+k_j-\rho_{ij}^2)$,
  \[
    \ip{W-U}{C}=\tfrac12\sum_j(c_j-1)k_j-\tfrac12\sum_{ij}(W_{ij}-U_{ij})\rho_{ij}^2,
  \]
  with $c_j=\sum_iW_{ij}$: a hub term in the communalities $k_j$ plus a
  proximity term in the exposure distances $\rho_{ij}$.
\end{theorem}
\begin{proof}\leanok
  Row sums of $W$ and $U$ are one and column sums of $U$ are one; collect terms.
\end{proof}

\begin{lemma}[Placebos keep the hub term]\label{thm:value-placebo}
  \lean{FreeFoundations.Value.placebo_hub}\leanok
  \uses{thm:value-decomposition}
  When communality is constant within the groups a permutation respects, the
  hub term is unchanged.
\end{lemma}
\begin{proof}\leanok
  Reindex the sum by the permutation.
\end{proof}

% ---------------------------------------------------------------------------
\section{The spatial likelihood}\label{sec:bp-likelihood}

Write $\Sigma=A^{-1}DA^{-\top}$ with $A=I-\lambda W$ and $D=\diag(d)$, and let
$M$ be the second moment of the residuals.

\begin{lemma}[Precision]\label{thm:likelihood-precision}
  \lean{FreeFoundations.Likelihood.precision_eq, FreeFoundations.Likelihood.trace_precision}\leanok
  $\Sigma^{-1}=A^\top D^{-1}A$, and
  $\tr(A^\top\diag(e)AM)=\sum_ie_i(AMA^\top)_{ii}$.
\end{lemma}
\begin{proof}\leanok
  Invert the product; cycle the trace.
\end{proof}

\begin{lemma}[Quadratic and log-determinant terms]\label{thm:likelihood-terms}
  \lean{FreeFoundations.Likelihood.quadratic_term, FreeFoundations.Likelihood.logdet_eq}\leanok
  $\tr(\Sigma^{-1}M)=\sum_im_i(\lambda)/d_i$ with filtered second moments
  $m_i(\lambda)=(AMA^\top)_{ii}$, and
  $\log\det\Sigma=\sum_i\log d_i-2\log\abs{\det A}$.
\end{lemma}
\begin{proof}\leanok
  \uses{thm:likelihood-precision}
  Lemma~\ref{thm:likelihood-precision}, and multiplicativity of the
  determinant.
\end{proof}

\begin{theorem}[Expected loss of the spatial-error forecast]\label{thm:likelihood-loss}
  \lean{FreeFoundations.Likelihood.spatial_loss}\leanok
  \uses{def:score-loss}
  \[
    L(\Sigma,M)=\sum_i\log d_i-2\log\abs{\det A}+\sum_i\frac{m_i(\lambda)}{d_i}:
  \]
  variance, spectrum and the $D^{-1}$-weighted filtered second moments.
\end{theorem}
\begin{proof}\leanok
  \uses{thm:likelihood-terms}
  Add the two terms of Lemma~\ref{thm:likelihood-terms}.
\end{proof}

\begin{definition}[Common-plus-idiosyncratic second moment]\label{def:likelihood-moment}
  \lean{FreeFoundations.Likelihood.moment}\leanok
  $M=\bigl(\ip{a_i}{a_j}\bigr)_{ij}+\diag(\eta)$: exposures $a_i$ to orthonormal
  common shocks plus idiosyncratic variances $\eta_i$.
\end{definition}

\begin{lemma}[Quadratic forms of the two parts]\label{thm:likelihood-parts}
  \lean{FreeFoundations.Likelihood.gram_quadratic, FreeFoundations.Likelihood.diagonal_quadratic}\leanok
  $\sum_{jk}x_j\ip{a_j}{a_k}x_k=\norm{\sum_jx_ja_j}^2$ and
  $\sum_{jk}x_j\eta_j\mathbb 1_{j=k}x_k=\sum_jx_j^2\eta_j$.
\end{lemma}
\begin{proof}\leanok
  Bilinearity of the inner product; the diagonal sum collapses.
\end{proof}

\begin{theorem}[Filtered second moment]\label{thm:likelihood-filtered}
  \lean{FreeFoundations.Likelihood.filtered_second_moment}\leanok
  \uses{def:likelihood-moment}
  For a zero-diagonal network,
  \[
    m_i(\lambda)=\Bigl\lVert a_i-\lambda\sum_jW_{ij}a_j\Bigr\rVert^2+\eta_i
    +\lambda^2\sum_jW_{ij}^2\eta_j :
  \]
  the neighbours' reconstruction error of firm $i$'s exposure, plus its own
  noise, plus the noise its neighbours bring in.
\end{theorem}
\begin{proof}\leanok
  \uses{thm:likelihood-parts}
  Row $i$ of $I-\lambda W$ is $e_i-\lambda W_{i\cdot}$; apply
  Lemma~\ref{thm:likelihood-parts} to it, using $W_{ii}=0$.
\end{proof}

\begin{corollary}[In powers of $\lambda$]\label{thm:likelihood-expand}
  \lean{FreeFoundations.Likelihood.filtered_second_moment_expand}\leanok
  \uses{def:likelihood-moment}
  \[
    m_i(\lambda)=(\norm{a_i}^2+\eta_i)-2\lambda\sum_jW_{ij}\ip{a_i}{a_j}
    +\lambda^2\Bigl(\Bigl\lVert\sum_jW_{ij}a_j\Bigr\rVert^2+\sum_jW_{ij}^2\eta_j\Bigr).
  \]
  Only common covariance enters the linear term; with equal $d_i$ it sums to
  $-2\ip{W}{C}/d$, the inner product of Theorem~\ref{thm:value-decomposition}.
\end{corollary}
\begin{proof}\leanok
  \uses{thm:likelihood-filtered}
  Expand the square in Theorem~\ref{thm:likelihood-filtered}.
\end{proof}

\begin{lemma}[Concentrated rows bring in more noise]\label{thm:likelihood-herfindahl}
  \lean{FreeFoundations.Likelihood.herfindahl_ge, FreeFoundations.Likelihood.herfindahl_uniform}\leanok
  A weight vector summing to one over $m$ firms has Herfindahl index
  $\sum_jw_j^2\ge1/m$, with equality for equal weights. The noise cost
  $\sum_jW_{ij}^2\eta_j$ is $\eta$ times this index when noise is equal.
\end{lemma}
\begin{proof}\leanok
  Cauchy--Schwarz: $1=(\sum_jw_j)^2\le m\sum_jw_j^2$.
\end{proof}

% ---------------------------------------------------------------------------
\section{Financial consequences}\label{sec:bp-stability}

\begin{lemma}[Inverse difference]\label{thm:stability-inverse}
  \lean{FreeFoundations.Stability.inverse_sub, FreeFoundations.Stability.norm_inverse_sub_le}\leanok
  For a left inverse $P$ of $X$ and a right inverse $P'$ of $X'$ in a normed
  ring, $P-P'=P(X'-X)P'$, so $\norm{P-P'}\le\norm P\norm{X'-X}\norm{P'}$.
\end{lemma}
\begin{proof}\leanok
  $P-P'=PX'P'-PXP'$.
\end{proof}

\begin{lemma}[Resolvent]\label{thm:stability-resolvent}
  \lean{FreeFoundations.Stability.resolvent_sub_le}\leanok
  If both resolvents have norm at most $M$,
  $\norm{(I-\lambda W)^{-1}-(I-\lambda W')^{-1}}\le\abs\lambda M^2\norm{W-W'}$.
\end{lemma}
\begin{proof}\leanok
  \uses{thm:stability-inverse}
  Lemma~\ref{thm:stability-inverse} with $X'-X=\lambda(W-W')$.
\end{proof}

\begin{theorem}[Covariance stability]\label{thm:stability-covariance}
  \lean{FreeFoundations.Stability.covariance_sub_le, FreeFoundations.Stability.matrix_covariance_sub_le}\leanok
  With common $\lambda$ and innovation covariance $V$, and resolvents of norm
  at most $M$,
  \[
    \norm{PVP^\top-P'VP'^\top}_2\le2\norm V_2\,\abs\lambda M^3\norm{W-W'}_2 ,
  \]
  for $P=(I-\lambda W)^{-1}$, $P'=(I-\lambda W')^{-1}$; in the spectral norm for
  real matrices.
\end{theorem}
\begin{proof}\leanok
  \uses{thm:stability-resolvent}
  Write the difference as $(P-P')VP^\top+P'V(P-P')^\top$ and apply
  Lemma~\ref{thm:stability-resolvent}.
\end{proof}

\begin{lemma}[Loss, quadratic term]\label{thm:stability-quadratic}
  \lean{FreeFoundations.Stability.abs_quadratic_le, FreeFoundations.Stability.quadratic_inverse_sub_le}\leanok
  $\abs{\ip{Tv}{v}}\le\norm T\norm v^2$, and for inverses $Q=\Sigma^{-1}$,
  $\tilde Q=\tilde\Sigma^{-1}$,
  \[
    \abs{u'Qu-u'\tilde Qu}\le\norm Q\norm{\tilde Q}\norm{\Sigma-\tilde\Sigma}\norm u^2 .
  \]
\end{lemma}
\begin{proof}\leanok
  \uses{thm:stability-inverse}
  Cauchy--Schwarz, then Lemma~\ref{thm:stability-inverse}.
\end{proof}

\begin{theorem}[Portfolio regret]\label{thm:stability-portfolio}
  \lean{FreeFoundations.Stability.portfolio_regret, FreeFoundations.Stability.long_only_sq_le}\leanok
  If $\hat w$ minimizes estimated variance over a feasible set whose portfolios
  have $\norm w\le B$, its true variance exceeds that of any feasible portfolio
  by at most $2B^2\norm{\hat\Sigma-\Sigma}$. Long-only, fully invested
  portfolios have $B=1$.
\end{theorem}
\begin{proof}\leanok
  \uses{thm:stability-quadratic}
  Add and subtract the estimated variances and bound each gap by
  Lemma~\ref{thm:stability-quadratic}; $\sum_iw_i^2\le(\sum_iw_i)^2=1$ for
  $w\ge0$.
\end{proof}


\begin{thebibliography}{9}
  \bibitem{fernandez}
  Fern\'andez, V. (2011).
  Spatial linkages in international financial markets.
  \emph{Quantitative Finance}, 11(2), 237--245.
  \url{https://doi.org/10.1080/14697680903127403}.
  \bibitem{assetembeddings}
  Gabaix, X., Koijen, R. S. J., Richmond, R. J., and Yogo, M. (2025).
  Asset Embeddings. BFI Working Paper 2025-51, April version. % chktex 8
  % chktex 8
  \bibitem{fulleotto}
  F\"ulle, M. J., and Otto, P. (2024).
  Spatial GARCH models for unknown spatial locations---an application to
  financial stock returns.\ \emph{Spatial Economic Analysis}, 19(1), 92--105.
  \url{https://doi.org/10.1080/17421772.2023.2237067}.
  \bibitem{qulee}
  Qu, X., and Lee, L.-F. (2015).
  Estimating a spatial autoregressive model with an endogenous
  spatial weight matrix.
  \emph{Journal of Econometrics}, 184(2), 209--232.
  \url{https://doi.org/10.1016/j.jeconom.2014.08.008}.
  \bibitem{pfn}
  M\"uller, S., Hollmann, N., Pineda Arango, S., Grabocka, J., and Hutter, F.
  (2022). Transformers can do Bayesian inference.\ \emph{International
  Conference on Learning Representations}.
  \bibitem{nagler}
  Nagler, T. (2023). Statistical foundations of prior-data fitted networks.
  \emph{Proceedings of the 40th International Conference on Machine Learning},
  PMLR 202.
  \bibitem{gneiting}
  Gneiting, T. (2011). Making and evaluating point forecasts.
  \emph{Journal of the American Statistical Association}, 106(494), 746--762.
  \url{https://doi.org/10.1198/jasa.2011.r10138}.
  \bibitem{secapi}
  US Securities and Exchange Commission.
  EDGAR Application Programming Interfaces.
  \url{https://www.sec.gov/search-filings/edgar-application-programming-interfaces}.
\end{thebibliography}
\end{document}
